\documentclass[10pt, a4paper]{article}

\usepackage[T1]{fontenc}
\usepackage[utf8]{inputenc}
\usepackage{tgpagella}
\usepackage{microtype}
\usepackage[spanish, es-noshorthands]{babel}
\usepackage{amsmath, amssymb, bm}
\usepackage{multicol}
\usepackage[most]{tcolorbox}
\usepackage[margin=1.5cm]{geometry}
\usepackage{fancyhdr}

\pagestyle{fancy}
\fancyhf{}
\setlength{\headheight}{16pt}
\lhead{\textbf{IMT-342} | Robótica}
\chead{Formulario Oficial — EC2}
\rhead{Unidad II}
\lfoot{Material Autorizado}
\rfoot{Página \thepage}
\renewcommand{\headrulewidth}{0.4pt}
\renewcommand{\footrulewidth}{0.4pt}

\definecolor{ucbblue}{RGB}{0, 51, 102}

% Estilo para las secciones del formulario
\tcbset{
    sectionbox/.style={
        colback=ucbblue!5,
        colframe=ucbblue,
        fonttitle=\bfseries,
        left=2mm, right=2mm, top=1mm, bottom=1mm,
        boxsep=0mm,
        arc=1mm
    }
}

\begin{document}

\begin{center}
    \Large\textbf{\textcolor{ucbblue}{Formulario Oficial — Evaluación EC2}}\\[0.2cm]
\end{center}

\begin{tcolorbox}[colback=yellow!10, colframe=black, title=\textbf{Reglas de Uso durante la Evaluación}, left=2mm, right=2mm, top=1mm, bottom=1mm]
    \small
    \textbf{Uso permitido en EC2:} una copia de este formulario + calculadora científica no programable[cite: 10]. \\
    \textbf{No permitido:} apuntes adicionales, CAS, laptop, Python/Jupyter, material propio[cite: 10].
\end{tcolorbox}

\vspace{0.2cm}
\begin{multicols*}{2}
\small

% --- SECCIÓN 1 ---
\begin{tcolorbox}[sectionbox, title=1. Cinemática Inversa / Pieper{[cite: 10]}]
\textbf{Pose homogénea}[cite: 10]
\begin{equation*}
    {}^0T_6^d=
    \begin{bmatrix}
    {}^0R_6^d & {}^0p_6^d\\
    0&1
    \end{bmatrix}
\end{equation*}

\textbf{Centro de muñeca (Modelo de clase)}[cite: 10]
\begin{equation*}
    \boxed{ {}^0p_W = {}^0p_6^d - d_6\,{}^0z_6^d }
\end{equation*}
con:
\begin{equation*}
    \boxed{ {}^0z_6^d = {}^0R_6^d \begin{bmatrix} 0\\0\\1 \end{bmatrix} }
\end{equation*}

\textbf{Separación brazo / muñeca}[cite: 10]
\begin{equation*}
    p_W \rightarrow (q_1, q_2, q_3)
\end{equation*}
\begin{equation*}
    \boxed{ {}^3R_6 = ({}^0R_3)^T{}^0R_6^d }
\end{equation*}
\begin{equation*}
    {}^3R_6 \rightarrow (q_4, q_5, q_6)
\end{equation*}

\textbf{Recuperación angular}[cite: 10]
\begin{equation*}
    \boxed{ q=\operatorname{atan2}(s,c) }
\end{equation*}

\textbf{Verificación}[cite: 10]
\begin{equation*}
    \boxed{ {}^0T_6(\mathbf{q}^{(k)}) \approx {}^0T_6^d }
\end{equation*}
\end{tcolorbox}

% --- SECCIÓN 2 ---
\begin{tcolorbox}[sectionbox, title=2. Jacobiano Geométrico{[cite: 10]}]
\textbf{Relación cinemática diferencial}[cite: 10]
\begin{equation*}
    \boxed{ v_e= \begin{bmatrix} v_L\\ \omega \end{bmatrix} = J(\mathbf{q})\dot{\mathbf{q}} }
\end{equation*}

\textbf{Columna de articulación revoluta}[cite: 10]
\begin{equation*}
    \boxed{ J_i= \begin{bmatrix} z_{i-1}\times(o_n-o_{i-1})\\ z_{i-1} \end{bmatrix} }
\end{equation*}

\textbf{Columna de articulación prismática}[cite: 10]
\begin{equation*}
    \boxed{ J_i= \begin{bmatrix} z_{i-1}\\ 0 \end{bmatrix} }
\end{equation*}

\textit{\textbf{Recordatorio:} Todos los vectores utilizados en una misma columna deben estar expresados en el \textbf{mismo sistema de referencia}[cite: 10].}
\end{tcolorbox}

% --- SECCIÓN 3 ---
\begin{tcolorbox}[sectionbox, title=3. Singularidades{[cite: 10]}]
\textbf{Definición general}[cite: 10]
\begin{equation*}
    \boxed{ \operatorname{rank}J(\mathbf{q}) < \operatorname{rank}_{\max} }
\end{equation*}
con $\operatorname{rank}_{\max} = \min(m,n)$ para $J\in\mathbb{R}^{m\times n}$[cite: 10].

\textbf{Caso cuadrado}[cite: 10] \\
Si $J$ es cuadrado, $\boxed{ \det J=0 }$ es un criterio equivalente de pérdida de rango completo[cite: 10].

\textbf{2R planar — referencia}[cite: 10]
\begin{equation*}
    \boxed{ \det J_p = a_1a_2\sin q_2 }
\end{equation*}
\end{tcolorbox}

% --- SECCIÓN 4 ---
\begin{tcolorbox}[sectionbox, title=4. Trayectoria LSPB{[cite: 10]}]
\textbf{Variables:} $q_i, q_f, t_f, \ddot{q}_c, t_c$, con $\dot{q}(0)=0$ y $\dot{q}(t_f)=0$[cite: 10].

\textbf{Aceleración mínima}[cite: 10]
\begin{equation*}
    \boxed{ |\ddot{q}_c| \ge \frac{4|q_f-q_i|}{t_f^2} }
\end{equation*}

\textbf{Tiempo de blend (signos consistentes)}[cite: 10]
\begin{equation*}
    \boxed{ t_c = \frac{t_f}{2} - \frac{1}{2} \sqrt{ t_f^2 - \frac{4(q_f-q_i)}{\ddot{q}_c} } }
\end{equation*}

\textbf{Velocidad de crucero}[cite: 10]
\begin{equation*}
    \boxed{ \dot{q}_c=\ddot{q}_c t_c }
\end{equation*}

\textbf{Fase 1 — aceleración ($0\le t\le t_c$)}[cite: 10]
\begin{equation*}
    \boxed{ q(t)=q_i+\frac{1}{2}\ddot{q}_c t^2 } \quad \boxed{ \dot{q}(t)=\ddot{q}_c t } \quad \boxed{ \ddot{q}(t)=\ddot{q}_c }
\end{equation*}

\textbf{Fase 2 — vel. constante ($t_c<t\le t_f-t_c$)}[cite: 10]
\begin{equation*}
    \boxed{ q(t) = q_i+ \ddot{q}_c t_c \left( t-\frac{t_c}{2} \right) }
\end{equation*}
\begin{equation*}
    \boxed{ \dot{q}(t)=\dot{q}_c } \quad \boxed{ \ddot{q}(t)=0 }
\end{equation*}

\textbf{Fase 3 — deceleración ($t_f-t_c<t\le t_f$)}[cite: 10]
\begin{equation*}
    \boxed{ q(t) = q_f - \frac{1}{2}\ddot{q}_c(t_f-t)^2 }
\end{equation*}
\begin{equation*}
    \boxed{ \dot{q}(t) = \ddot{q}_c(t_f-t) } \quad \boxed{ \ddot{q}(t) = -\ddot{q}_c }
\end{equation*}
\end{tcolorbox}

% --- SECCIÓN 5 ---
\begin{tcolorbox}[sectionbox, title=5. Dinámica — Euler–Lagrange{[cite: 10]}]
\textbf{Lagrangiano y Ecuación General}[cite: 10]
\begin{equation*}
    \boxed{ L=T-V } \quad \boxed{ \frac{d}{dt} \left( \frac{\partial L}{\partial\dot{q}_i} \right) - \frac{\partial L}{\partial q_i} = \tau_i }
\end{equation*}

\textbf{Modelo 1R trabajado}[cite: 10]
\begin{align*}
    \boxed{ J=I_c+m\ell_c^2 } \quad &\boxed{ T= \frac{1}{2}J\dot{q}^2 } \\
    \boxed{ V= mg\ell_c(1-\cos q) } \quad &\boxed{ \tau = J\ddot{q} + mg\ell_c\sin q }
\end{align*}

\textbf{Forma general del manipulador}[cite: 10]
\begin{equation*}
    \boxed{ \tau = M(\mathbf{q})\ddot{\mathbf{q}} + C(\mathbf{q},\dot{\mathbf{q}})\dot{\mathbf{q}} + \mathbf{g}(\mathbf{q}) }
\end{equation*}
\end{tcolorbox}

% --- SECCIÓN 6 ---
\begin{tcolorbox}[sectionbox, title=6. Mapa Conceptual de Unidad II{[cite: 10]}]
\begin{align*}
    \boxed{ T_d \rightarrow \text{IK / Pieper} \rightarrow \mathbf{q} } \\
    \boxed{ q_i,q_f,\text{ tiempo} \rightarrow \text{LSPB} \rightarrow q,\dot{q},\ddot{q} } \\
    \boxed{ q,\dot{q} \rightarrow J(\mathbf{q}) \rightarrow \dot{\mathbf{x}} } \\
    \boxed{ J(\mathbf{q}) \rightarrow \operatorname{rank}J \rightarrow \text{singularidad} } \\
    \boxed{ q,\dot{q},\ddot{q} \rightarrow \text{dinámica} \rightarrow \tau }
\end{align*}
\end{tcolorbox}

% --- SECCIÓN 8 ---
\begin{tcolorbox}[sectionbox, title=8. Identidades Útiles{[cite: 10]}]
\begin{equation*}
    \sin^2\theta+\cos^2\theta=1
\end{equation*}
\begin{equation*}
    \sin(A-B) = \sin A\cos B-\cos A\sin B
\end{equation*}
\begin{equation*}
    R^{-1}=R^T \quad \text{(Para matriz de rotación válida)}
\end{equation*}
\begin{equation*}
    a\times b = \begin{bmatrix} a_yb_z-a_zb_y\\ a_zb_x-a_xb_z\\ a_xb_y-a_yb_x \end{bmatrix}
\end{equation*}
\end{tcolorbox}

% --- SECCIÓN 7 ---
\begin{tcolorbox}[sectionbox, title=7. Recordatorios de Interpretación{[cite: 10]}]
\begin{itemize}
    \item Una rama de IK debe \textbf{verificarse} con FK[cite: 10].
    \item \texttt{atan2} preserva información de cuadrante[cite: 10].
    \item El Jacobiano \textbf{no} es una matriz homogénea[cite: 10].
    \item Una columna de $J$ representa una contribución de velocidad[cite: 10].
    \item Singularidad significa \textbf{pérdida de rango}, no simplemente “determinante cero”[cite: 10].
    \item Un punto en configuración no es una trayectoria[cite: 10].
    \item En LSPB: $q(t), \dot{q}(t)$ son continuos en el perfil ideal trabajado[cite: 10].
    \item Que $\ddot{q}=0$ no implica $\tau=0$[cite: 10].
    \item El modelo 1R tiene un término de velocidad nulo en este caso particular; no debe generalizarse a todos los manipuladores[cite: 10].
\end{itemize}
\end{tcolorbox}

% --- SECCIÓN 9 ---
\begin{tcolorbox}[colback=white, colframe=black, title=\textbf{9. Lo que este formulario NO contiene}]
Este formulario no incluye: soluciones de problemas, ramas de Pieper completas, Jacobianos resueltos, singularidades del ABB IRB120, parámetros dinámicos del ABB, ejercicios resueltos, algoritmos de pseudoinversa, ni redundancia/espacio nulo[cite: 10]. La intención es recordar la \textbf{estructura matemática}, no resolver el examen por el estudiante[cite: 10].
\end{tcolorbox}

\end{multicols*}

\end{document}
